\documentclass[../../main.tex]{subfiles}

\begin{document}

\chapter{Interupsi BIOS/DOS}
\label{chap:interupsi}

% =========================================================
% OUTCOME PEMBELAJARAN KHUSUS BAB
% =========================================================

\begin{subcpmkbox}
\vspace{1.1em}
\SubCPMK{2.3} --- Menjelaskan mekanisme interupsi dan tabel vektor
interupsi serta menggunakan layanan BIOS/DOS (INT 21h, INT 10h,
INT 16h) untuk operasi masukan dan keluaran.

\vspace{0.4em}
\textbf{CPMK yang didukung:} \CPMK{02} \quad\textbullet\quad
\textbf{CPL yang didukung:} \CPL{02}
\end{subcpmkbox}

\begin{tujuanbox}
\vspace{1.1em}
Setelah mempelajari bab ini, mahasiswa mampu:
\begin{enumerate}
    \item menjelaskan pengertian interupsi perangkat keras dan
          perangkat lunak;
    \item menjelaskan tabel vektor interupsi dan cara menentukan
          alamat handler;
    \item menjelaskan penyimpanan dan pemulihan keadaan prosesor
          pada interupsi;
    \item menggunakan layanan video BIOS \code{INT 10h};
    \item menggunakan layanan \eng{keyboard} BIOS \code{INT 16h};
    \item menggunakan layanan DOS \code{INT 21h} untuk mencetak
          karakter, mencetak string, membaca masukan, dan mengakhiri
          program;
    \item membedakan \code{INT 20h} dan \code{INT 21h} fungsi
          \hex{4C}.
\end{enumerate}
\end{tujuanbox}

\begin{materibox}
\vspace{1.1em}
\begin{enumerate}
    \item Pengertian interupsi perangkat keras dan perangkat lunak.
    \item Tabel vektor interupsi 1 KB dan susunan vektornya.
    \item Mekanisme \code{INT} dan \code{IRET}.
    \item BIOS interrupt: \code{INT 10h} dan \code{INT 16h}.
    \item DOS interrupt: \code{INT 21h} fungsi \hex{01}, \hex{02},
          \hex{09}, \hex{0A}, dan \hex{4C}.
    \item \code{INT 20h} dan pengakhiran program.
    \item Operasi masukan dan keluaran dengan BIOS/DOS.
\end{enumerate}
\end{materibox}

\subfile{section-01}
\subfile{section-02}
\subfile{section-03}

\section{Contoh Program}
\label{sec:contoh-bab9}

Bagian ini menyajikan dua program utuh. Program pertama memakai
seluruh direktif data, sedangkan program kedua memakai \dir{EQU}
sebagai konstanta sekaligus memperlihatkan apa yang terjadi bila
direktif itu didefinisikan ulang.

\subsection{Contoh 1: Mendemonstrasikan Semua Direktif Data}

Program pada Kode~\ref{lst:direktif-data} memakai \dir{DB},
\dir{DW}, \dir{DD}, \dir{DUP}, dan penanda \code{?}, lalu mencetak
sebagian isi memori tersebut ke layar.

\lstinputlisting[
    language=MASM,
    caption={Program \file{direktif-data.asm}: direktif data dan
    isi memorinya.},
    label={lst:direktif-data}
]{\CodeDir/direktif-data.asm}

Bagian \dir{.DATA} pada program tersebut dapat dibaca sebagai
berikut.

\begin{itemize}
    \item \lstasm|sapaan DB 'Halo, assembly!', 0Dh, 0Ah, '$'|
          menyediakan \dec{18} byte: \dec{15} byte isi string,
          satu byte \code{0Dh} dan satu byte \code{0Ah} untuk
          pindah baris, serta \code{'\$'} sebagai penanda akhir
          untuk layanan DOS fungsi \hex{09}.
    \item \lstasm|panjang EQU $ - sapaan| meminta assembler
          menghitung panjang itu sendiri. Nilainya terlihat pada
          berkas \file{.lst} (Kode~\ref{lst:listing-lst}).
    \item \lstasm|kata DW 4241h, 4443h| menyediakan empat byte.
          Karena urutan \eng{little-endian}, byte yang tersimpan
          berurutan \hex{41}, \hex{42}, \hex{43}, \hex{44}.
    \item \lstasm|nilai DD 12345678h| menyediakan empat byte tanpa
          dipakai program, tetapi tetap menggeser \eng{offset}
          \code{bintang} dan \code{buffer}.
    \item \lstasm|bintang DB 4 DUP('*')| menunjukkan \dir{DUP} yang
          berisi nilai awal, sedangkan \lstasm|buffer DB 8 DUP(?)|
          hanya menyediakan ruang kosong.
\end{itemize}

Bagian \dir{.CODE} mencetak sapaan, lalu membaca \code{kata}
\emph{byte per byte} dengan \lstasm|BYTE PTR [kata+1]|. Cara
pengalamatan bertanda kurung siku dibahas pada
Bab~\ref{chap:mode-pengalamatan}.

\begin{session}[caption={Keluaran program \file{direktif-data.asm}.}, label={lst:keluaran-direktif-data}]
Halo, assembly!
ABCD****
\end{session}

Baris kedua keluaran tersebut membuktikan urutan
\eng{little-endian}: data \hex{4241} dan \hex{4443} terbaca di
memori sebagai \code{ABCD}.

\subsection{Contoh 2: Konstanta \dir{EQU}}

Program pada Kode~\ref{lst:equ-konstanta} mengganti setiap nomor
fungsi DOS dengan nama yang bermakna. Nilainya tidak berubah
dibandingkan penulisan langsung; yang berubah hanyalah keterbacaan
program.

\lstinputlisting[
    language=MASM,
    caption={Program \file{equ-konstanta.asm}: konstanta simbolik
    dengan \dir{EQU}.},
    label={lst:equ-konstanta}
]{\CodeDir/equ-konstanta.asm}

\begin{session}[caption={Keluaran program \file{equ-konstanta.asm}.}, label={lst:keluaran-equ-konstanta}]
Konstanta EQU
X = 7
\end{session}

\subsection{Definisi Ulang \dir{EQU} Dibandingkan dengan \code{=}}

Kode~\ref{lst:equ-ulang-gagal} menimbulkan kesalahan perakitan yang
sama seperti pada Kode~\ref{lst:equ-salah}.

\begin{lstlisting}[
    language=MASM,
    caption={Baris yang menimbulkan kesalahan perakitan.},
    label={lst:equ-ulang-gagal}
]
NILAI   EQU 5           ; definisi pertama
NILAI   EQU 9           ; gagal: Symbol already defined
\end{lstlisting}

Bila kata kunci \dir{EQU} diganti dengan \code{=}, kedua baris itu
menjadi sah. Program berikut mencetak dua angka berbeda dari satu
nama simbol.

\begin{lstlisting}[
    language=MASM,
    caption={\code{=} boleh didefinisikan ulang; program mencetak
    \code{59}.},
    label={lst:sama-dengan-ulang}
]
.MODEL SMALL
.STACK 100h

.DATA
    awalan  DB 'Nilai = ', '$'

.CODE
START:
    MOV AX, @DATA
    MOV DS, AX

    MOV AH, 09h             ; cetak awalan
    LEA DX, awalan
    INT 21h

NILAI   = 5                 ; definisi pertama
    MOV AH, 02h             ; cetak hasil '0' + 5 = '5'
    MOV DL, '0' + NILAI
    INT 21h

NILAI   = 9                 ; definisi ulang: sah
    MOV AH, 02h             ; cetak hasil '0' + 9 = '9'
    MOV DL, '0' + NILAI
    INT 21h

    MOV AH, 4Ch
    MOV AL, 00h
    INT 21h

END START
\end{lstlisting}

\begin{session}[caption={Keluaran program pada Kode~\ref{lst:sama-dengan-ulang}.}]
Nilai = 59
\end{session}

Kesimpulannya sederhana: pakai \dir{EQU} untuk nilai yang
\emph{tidak akan pernah berubah}, seperti nomor fungsi DOS dan
ukuran tetap, dan pakai \code{=} untuk nilai yang memang harus
diganti di tengah berkas.

\section{Praktikum}
\label{sec:praktikum-bab9}

\begin{praktikumbox}
\vspace{1.1em}
\textbf{Judul:} Menggunakan layanan BIOS/DOS dan memeriksa tabel
vektor interupsi.\\[0.3em]
\textbf{Sub-CPMK:} 2.3 \quad\textbullet\quad
\textbf{Alokasi waktu:} 100 menit \quad\textbullet\quad
\textbf{Bentuk:} individu
\end{praktikumbox}

\subsection{Tujuan Praktikum}

Setelah menyelesaikan praktikum ini, mahasiswa mampu menjalankan
program dengan layanan \code{INT 21h}, \code{INT 10h}, dan
\code{INT 16h}, serta menjelaskan isi tabel vektor interupsi.

\subsection{Alat dan Bahan}

\begin{itemize}
    \item Komputer dengan emu8086 terpasang.
    \item DOSBox dengan TASM, TLINK, dan DEBUG.
    \item Berkas \file{input-nama.asm} dan \file{warna-layar.asm}
          dari \file{code/asm/}.
    \item Lembar jurnal praktikum.
\end{itemize}

\subsection{Langkah Kerja}

\begin{enumerate}
    \item \textbf{Menguji masukan dan keluaran DOS.}
    \begin{enumerate}
        \item Rakit \file{input-nama.asm} dengan TASM, lalu link
              dengan TLINK.
        \item Jalankan program dan catat pertanyaan serta keluaran
              yang muncul.
        \item Buka program dengan DEBUG, lalu periksa penyangga
              \code{nama} pada segmen data.
        \item Catat byte kapasitas, byte jumlah karakter, dan
              karakter yang tersimpan.
    \end{enumerate}

    \item \textbf{Menguji layanan video BIOS.}
    \begin{enumerate}
        \item Rakit dan jalankan \file{warna-layar.asm}.
        \item Catat posisi judul dan garis yang tampil.
        \item Ubah atribut warna judul dari \hex{0E} menjadi
              \hex{0A}, lalu rakit ulang dan amati perubahannya.
        \item Ubah posisi baris judul dari 5 menjadi 3 dan kolom dari
              10 menjadi 20. Catat hasilnya.
    \end{enumerate}

    \item \textbf{Memeriksa tabel vektor interupsi.}
    \begin{enumerate}
        \item Jalankan DEBUG.
        \item Tampilkan vektor \code{INT 21h} dengan perintah
              \lstasm|D 0000:0084 L 4|.
        \item Catat dua byte offset dan dua byte segmen yang tampil.
        \item Hitung alamat fisik handler dengan rumus
              $(\text{segmen} \times \hex{10}) + \text{offset}$.
        \item Ulangi untuk \code{INT 10h} pada alamat
              \hex{0000}:\hex{0040}.
    \end{enumerate}

    \item \textbf{Menguji interupsi keyboard BIOS.}
    \begin{enumerate}
        \item Tulis program berikut sebagai \file{baca-tombol.asm}.

\begin{lstlisting}[language=MASM, numbers=none]
.MODEL SMALL
.CODE
START:
    MOV AH, 00h     ; baca tombol, menunggu
    INT 16h         ; hasil di AH (scan) dan AL (ASCII)
    MOV AH, 4Ch
    INT 21h
END START
\end{lstlisting}

        \item Rakit, link, lalu jalankan program.
        \item Tekan beberapa tombol dan jelaskan mengapa program baru
              berakhir setelah tombol ditekan.
        \item Periksa isi \reg{AX} setelah \code{INT 16h} dengan
              \eng{debugger}.
    \end{enumerate}

    \item \textbf{Modifikasi program.}
    \begin{enumerate}
        \item Ubah pertanyaan pada \file{input-nama.asm} sesuai
              kebutuhan.
        \item Ubah teks dan warna pada \file{warna-layar.asm}.
        \item Rakit ulang dan pastikan perubahan terlihat.
    \end{enumerate}
\end{enumerate}

\subsection{Hasil yang Dikumpulkan}

\begin{enumerate}
    \item Berkas \file{input-nama.asm} dan \file{warna-layar.asm}
          versi modifikasi.
    \item Tangkapan layar keluaran program masukan.
    \item Tangkapan layar tampilan berwarna.
    \item Tabel vektor \code{INT 21h} dan \code{INT 10h} beserta
          alamat fisik handler.
    \item Berkas \file{baca-tombol.asm} dan penjelasan hasilnya.
\end{enumerate}

\begin{asesmenbox}
\vspace{1.1em}
\textbf{Asesmen praktikum 7 (bobot 5\% dari komponen praktikum).}

\begin{tabular}{@{}p{6.2cm}cp{4.6cm}@{}}
    \toprule
    \textbf{Indikator} & \textbf{Bobot} & \textbf{Kriteria} \\
    \midrule
    Program masukan/keluaran DOS berjalan & 30\% &
    Pertanyaan, masukan, dan sapaan tampil dengan benar \\
    \addlinespace
    Program layanan video berjalan & 25\% &
    Teks dan warna tampil sesuai register yang diisi \\
    \addlinespace
    Ketepatan membaca tabel vektor & 30\% &
    Offset, segmen, dan alamat fisik handler dihitung benar \\
    \addlinespace
    Analisis interupsi keyboard & 15\% &
    Perilaku menunggu dan isi register dijelaskan \\
    \bottomrule
\end{tabular}
\end{asesmenbox}

\section{Latihan}
\label{sec:latihan-bab9}

\begin{enumerate}
    \item \textbf{Konseptual.} Jelaskan perbedaan instruksi mesin dan
          direktif assembler, serta mengapa direktif disebut
          \eng{pseudo-op}. Berikan dua contoh direktif yang mengisi
          memori dan dua contoh yang tidak mempengaruhi isi berkas
          \file{.exe}.

    \item \textbf{Konseptual.} Sebutkan direktif data yang tepat
          untuk setiap kebutuhan berikut, lalu tuliskan
          deklarasinya: (a) satu nilai \hex{0A} untuk pindah baris;
          (b) sebuah tabel berisi delapan nilai 16-bit untuk
          \hex{0000}; (c) sebuah penampung 32 byte yang diisi saat
          program berjalan; (d) sebuah string \code{'SEMARANG'}
          yang siap dicetak dengan layanan DOS fungsi \hex{09}.

    \item \textbf{Analisis.} Diberikan blok data berikut. Hitung
          total byte yang disediakan dan tuliskan \eng{offset}
          setiap label.

\begin{lstlisting}[language=MASM, numbers=none]
.DATA
    a   DB 'TEST'
    b   DW 2 DUP(0)
    c   DD 1, 2
    d   DB 3 DUP(?)
    e   DT 0
\end{lstlisting}

    \item \textbf{Analisis \eng{little-endian}.} Tuliskan urutan byte
          di memori untuk deklarasi \lstasm|DW 0102h| dan
          \lstasm|DD 03040506h|. Jelaskan mengapa urutannya
          demikian, lalu tentukan nilai yang terbaca pada
          \reg{AX} bila program menjalankan \lstasm|MOV AX, w|
          dengan \lstasm|w DW 0102h|.

    \item \textbf{Analisis kesalahan.} Perhatikan potongan berikut.

\begin{lstlisting}[language=MASM, numbers=none]
BARIS   EQU 25
KOLOM   EQU 80
BARIS   EQU 50
\end{lstlisting}

          \begin{enumerate}
              \item Baris mana yang menimbulkan kesalahan, dan apa
                    pesan yang dikeluarkan TASM?
              \item Ubah potongan itu agar sah, tetapi tetap
                    menghasilkan \code{BARIS} bernilai \dec{50}
                    pada pemakaian terakhir.
              \item Jelaskan mengapa pemakaian \code{=} tidak
                    menimbulkan masalah pada kasus ini.
          \end{enumerate}

    \item \textbf{Analisis alur translasi.} Sebuah berkas sumber
          berakhir dengan \lstasm|END| tanpa nama label.
          \begin{enumerate}
              \item Apakah TASM melaporkan kesalahan? Jelaskan.
              \item Apakah TLINK melaporkan kesalahan? Jelaskan
                    pesan yang muncul.
              \item Bagaimana cara memperbaikinya?
          \end{enumerate}

    \item \textbf{Penulisan program.} Tulislah program lengkap yang
          mendeklarasikan sebuah tabel berisi lima nilai
          \hex{41}--\hex{45} dengan \dir{DB}, mencetaknya dengan
          layanan DOS fungsi \hex{02} satu karakter per
          pemanggilan, lalu keluar dengan fungsi \hex{4C}.

    \item \textbf{Penulisan program dan pembacaan berkas hasil.}
          Deklarasikan semua nomor fungsi DOS yang dipakai dengan
          \dir{EQU}, lalu rakit dengan \file{tasm /l} dan
          \file{tlink /m}. Salin tabel segmen dari berkas
          \file{.map}, lalu jelaskan hubungan besar segmen
          \file{\_DATA} dengan jumlah byte deklarasi data pada
          program Anda.
\end{enumerate}

\begin{refleksibox}
Jawab pertanyaan berikut pada jurnal praktikum Anda.
\begin{enumerate}
    \item Direktif mana yang paling sering Anda lupakan fungsinya?
          Bagaimana Anda mengingatnya kembali?
    \item Setelah membaca berkas \file{.lst} dan \file{.map},
          apakah gambaran Anda tentang kerja assembler dan linker
          menjadi lebih jelas? Bagian mana yang paling menolong?
    \item Kesalahan perakitan apa yang Anda temui pada latihan ini,
          dan apa pelajaran yang Anda ambil dari kesalahan
          tersebut?
\end{enumerate}
\end{refleksibox}

\section{Rangkuman}
\label{sec:rangkuman-bab9}

Pada bab ini telah dibahas:

\begin{itemize}
    \item Kode sumber assembly terdiri atas dua jenis baris:
          instruksi mesin yang menghasilkan opcode dan dieksekusi
          saat program berjalan, serta direktif assembler yang
          hanya berlaku pada saat perakitan dan tidak menghasilkan
          kode mesin \cite{hyde2010}.

    \item \dir{DB} menyediakan satu byte, \dir{DW} dua byte,
          \dir{DD} empat byte, dan \dir{DT} sepuluh byte. Operator
          \dir{DUP} mengulang satu pola, sedangkan penanda \code{?}
          hanya menyediakan ruang tanpa nilai awal.

    \item Bilangan multi-byte disimpan dengan urutan
          \eng{little-endian}, sehingga \lstasm|DW 1234h|
          tersimpan sebagai byte \hex{34} lalu \hex{12}.

    \item \dir{EQU} membuat konstanta tetap yang tidak dapat
          didefinisikan ulang, sedangkan \code{=} membuat nilai
          yang boleh diganti berkali-kali.

    \item Direktif lokasi dan struktur meliputi \dir{ORG},
          \dir{SEGMENT}/\dir{ENDS}, \dir{ASSUME}, dan \dir{END}.
          Direktif \dir{END} yang diberi label sekaligus
          menetapkan titik masuk program.

    \item \dir{.MODEL} memilih model memori: \file{TINY} satu
          segmen, \file{SMALL} satu segmen kode dan satu segmen
          data, sedangkan model yang lebih besar memakai banyak
          segmen.

    \item Assembler bekerja dalam dua lintasan: membangun tabel
          simbol dan menghitung \eng{offset}, lalu menghasilkan
          kode mesin. Berkas \file{.lst} memperlihatkan hasilnya,
          termasuk penanda \code{R} \cite{bodur}. Linker kemudian
          menyelesaikan simbol eksternal, melakukan
          \eng{relocation}, dan merekam alamat akhir pada berkas
          \file{.map} \cite{nasm2024}.

    \item Alur lengkap dari DOS adalah
          \file{tasm /l program.asm} $\rightarrow$
          \file{tlink /m program.obj} $\rightarrow$ menjalankan
          \file{program.exe}.
\end{itemize}

\begin{checklistbox}
Tandai setiap butir yang sudah Anda kuasai.

\begin{itemize}
    \item[\checkbox] Saya dapat menjelaskan mengapa direktif tidak
          menghasilkan kode mesin.
    \item[\checkbox] Saya dapat memilih direktif data yang tepat
          dan menghitung besar memori yang disediakannya.
    \item[\checkbox] Saya dapat menentukan urutan byte
          \eng{little-endian} pada \dir{DW} dan \dir{DD}.
    \item[\checkbox] Saya dapat menjelaskan perbedaan \dir{EQU} dan
          \code{=} beserta akibat definisi ulang.
    \item[\checkbox] Saya dapat menuliskan struktur segmen program
          \file{.COM} dan \file{.EXE}.
    \item[\checkbox] Saya dapat memilih model memori \dir{.MODEL}
          dan menjelaskan pengaruhnya pada jumlah segmen.
    \item[\checkbox] Saya dapat membaca tabel simbol pada berkas
          \file{.lst} dan tabel segmen pada berkas \file{.map}.
    \item[\checkbox] Saya dapat merakit dan mengaitkan program
          dengan TASM dan TLINK tanpa kesalahan.
\end{itemize}

\vspace{0.4em}
Bila masih ada butir yang belum bertanda, ulangi bagian yang
bersesuaian sebelum melanjutkan ke bab berikutnya.
\end{checklistbox}

\begin{infobox}
\textbf{Rujukan bab.} Pembahasan direktif data dan peran assembler
diperkuat oleh \cite{hyde2010} dan \cite{bodur}. Keterangan praktis
mengenai direktif \dir{EQU}, \dir{ORG}, dan \dir{DUP} pada assembler
modern dapat ditelusuri pada \cite{nasm2024}, sedangkan contoh
pemakaiannya pada lingkungan belajar tersedia pada
\cite{emu8086help}.
\end{infobox}

\section{Evaluasi}
\label{sec:evaluasi-bab9}

\subsection{Soal Pilihan Ganda}

\begin{enumerate}
    \item Pernyataan yang benar tentang direktif assembler adalah
          \dots
          \begin{enumerate}
              \item direktif menghasilkan opcode seperti instruksi
              \item direktif hanya berlaku saat perakitan dan tidak
                    menghasilkan kode mesin
              \item direktif dieksekusi prosesor saat program
                    berjalan
              \item direktif tidak berpengaruh pada berkas
                    \file{.exe}
          \end{enumerate}

    \item Direktif yang menyediakan dua byte per nilai adalah \dots
          \begin{enumerate}
              \item \dir{DB} \item \dir{DW}
              \item \dir{DD} \item \dir{DT}
          \end{enumerate}

    \item Deklarasi \lstasm|nilai DW 1234h| menempatkan byte di
          memori dengan urutan \dots
          \begin{enumerate}
              \item \hex{12} lalu \hex{34}
              \item \hex{34} lalu \hex{12}
              \item \hex{00} lalu \hex{12}
              \item \hex{43} lalu \hex{21}
          \end{enumerate}

    \item Pasangan direktif yang benar beserta sifatnya adalah
          \dots
          \begin{enumerate}
              \item \dir{EQU} dapat didefinisikan ulang, \code{=}
                    tidak dapat
              \item \dir{EQU} tidak dapat didefinisikan ulang,
                    \code{=} dapat
              \item keduanya tidak dapat didefinisikan ulang
              \item keduanya selalu dapat didefinisikan ulang
          \end{enumerate}

    \item Tugas yang dikerjakan linker, bukan assembler, adalah
          \dots
          \begin{enumerate}
              \item menghitung \eng{offset} label di dalam satu
                    segmen
              \item menerjemahkan mnemonic menjadi opcode
              \item menyelesaikan simbol eksternal dan melakukan
                    \eng{relocation}
              \item membaca direktif \dir{EQU}
          \end{enumerate}
\end{enumerate}

\subsection{Soal Uraian}

\begin{enumerate}
    \item Jelaskan perbedaan instruksi mesin dan direktif assembler,
          lalu berikan dua contoh direktif yang mempengaruhi isi
          berkas \file{.exe} dan dua contoh yang tidak.

    \item Tuliskan deklarasi data untuk: string \code{'PRAKTIKUM'},
          \eng{array} berisi empat nilai 32-bit \hex{00000001}, dan
          penampung \eng{byte} berisi enam belas nilai nol. Hitung
          total byte dan \eng{offset} setiap label.

    \item Jelaskan alur translasi \file{.asm} $\rightarrow$
          \file{.obj} $\rightarrow$ \file{.exe}. Sebutkan apa yang
          sudah diketahui assembler pada akhir lintasan pertama dan
          apa yang baru ditentukan linker. Kaitkan jawaban Anda
          dengan penanda \code{R} pada berkas \file{.lst}.

    \item Sebuah program dirakit tanpa pesan kesalahan, tetapi
          linker menolak berkas objeknya. Sebutkan tiga
          kemungkinan penyebabnya beserta cara memperbaikinya.
\end{enumerate}

\subsection{Kisi-kisi Penilaian}

\begin{table}[H]
    \centering
    \caption{Kisi-kisi evaluasi Bab~\ref{chap:translasi}.}
    \label{tab:kisi-bab9}
    \small
    \begin{tabular}{@{}p{1.5cm}p{5.0cm}p{2.6cm}c@{}}
        \toprule
        \textbf{Sub-CPMK} & \textbf{Indikator yang diukur} &
        \textbf{Bentuk soal} & \textbf{Bobot} \\
        \midrule
        2.3 & Membedakan instruksi dan direktif serta
        memilih direktif data & Pilihan ganda 1, 2 & 15\% \\
        \addlinespace
        2.3 & Menghitung ukuran dan \eng{offset} deklarasi data &
        Uraian 2 & 25\% \\
        \addlinespace
        2.3 & Menerapkan aturan \eng{little-endian} dan perbedaan
        \dir{EQU} dengan \code{=} & Pilihan ganda 3, 4; Uraian 1 &
        25\% \\
        \addlinespace
        2.3 & Menjelaskan tugas assembler dan linker serta membaca
        berkas hasil & Pilihan ganda 5; Uraian 3, 4 & 35\% \\
        \bottomrule
    \end{tabular}
\end{table}

\begin{checklistbox}
\textbf{Checklist Pencapaian Kompetensi Bab~\ref{chap:translasi}}

Seorang mahasiswa dinyatakan kompeten pada Sub-CPMK 2.3 apabila:
\begin{itemize}
    \item[\checked] dapat menjelaskan mengapa direktif tidak
          menghasilkan kode mesin;
    \item[\checked] dapat menuliskan deklarasi \dir{DB}, \dir{DW},
          \dir{DD}, dan \dir{DUP} beserta perhitungan ukurannya;
    \item[\checked] dapat menentukan urutan byte
          \eng{little-endian} pada data multi-byte;
    \item[\checked] dapat menjelaskan perbedaan \dir{EQU} dan
          \code{=} serta akibat definisi ulang;
    \item[\checked] dapat menjelaskan tugas assembler dan linker
          pada alur \file{.asm} $\rightarrow$ \file{.obj}
          $\rightarrow$ \file{.exe};
    \item[\checked] berhasil merakit dan mengaitkan program dengan
          TASM dan TLINK, serta membaca berkas \file{.lst} dan
          \file{.map}.
\end{itemize}
\end{checklistbox}


\end{document}
